\documentclass[class=NCU_thesis, crop=false, float=true]{standalone}

\begin{document}
% This file is common document region, 
% it will loading when LaTeX see ``\begin{document}'' 
% It's for some command which need use in document region.
% DON'T put preamble only command in here, eg. \usepackage
% --------------------------------------------------

% standalone setting for sub-tex. (not every class option can use \standaloneconfig)
\IfStandalone{\standaloneconfig{float=true}}{} 
% ``float=false'' disable floating environment when build sub-tex.

% Old method for change font size, but no limit for size so keep here.
% \fontsize{size}{baselineskip} 無級調大小、行高, 行高倍數 = baselineskip/size/1.2 ，詳見教學。
% \fontsize{14}{25}\selectfont  

\ifzh{
    \setlength{\parindent}{2em} % 中文段首縮排2中文字寬
}{}                             % 英文照 LaTeX 預設


\chapter{章名（章節示例）}
章內容
\section{節名}
節內容
\subsection{小節名}
 \subsubsection{小小節}
 \paragraph{段}


\chapter{文字}
第一行。
仍是第一行。 \\
第二行。


\chapter{圖片}
\subsection{插入單一圖片}
\insfig[0.15][fig:labal_test][!hbt]{logo-Linux.png}[caption][short caption]

\subsection{插入多張圖片}
\begin{figure}[!hbt]
    %\captionsetup[subfigure]{labelformat=empty} % 完全隱藏圖號
    \centering
    \subcaptionbox
        {Linux (A kernel of OS)
        \label{fig:subfig_linux}}
        {\includegraphics[width=0.3\linewidth]{logo-Linux.png}}
    ~
    \subcaptionbox
        {Debain (A popular distribution). Demo long caption,  bla bla bla bla bla.
        \label{fig:subfig_debian}}
        {\includegraphics[width=0.3\linewidth]{logo-Debian.png}}
    \vspace{\baselineskip} % 分隔上下列
    \subcaptionbox
        {GNU (A project of OS)
        \label{fig:subfig_gnu}}
        {\includegraphics[width=0.6\linewidth]{logo-GNU.png}}
     % use ``\subref{fig:subfig_debian}'' get ID of subfigure(this ID is Debian)
    \caption{caption, 使用 \subref{fig:subfig_debian}取得子圖（Debian）編號 }
    \label{fig:labal}
\end{figure}


\chapter{表格}
\subsection{一般表格}
\begin{table}[h]
    \centering
    \caption{Solution}
    \begin{tabular}{| l | l |}
        \hline
        Component & Concentration(mM) \\ \hline
        \ce{NaCl} & 118.0 \\ \hline
    \end{tabular}
\end{table}

\subsection{自動折行表格}
\begin{table}[h]
    \centering
    \begin{tabularx}{\textwidth}{| l | X |}
        \hline
        short & short short \\ \hline
        long & long long long long long long long long long long \\ \hline
    \end{tabularx}
\end{table}

\end{document}